\section{Trace, endpoint, and gap estimates}
\label{app:green-gap}

This appendix supplies the quantitative estimates suppressed in
\cref{thm:canonical-traces,prop:gap-expansion}.

\subsection{The nonlinear Green contraction}

Let \(d_0\) be the loss in \cref{lem:green}, choose \(c>c_0\) and
\(M>m_0+d_0+2\), and define
\[
 \|\xi\|_\cX
 =\sup_{J_\sigma}\e^{-c|s|}\langle s\rangle^{-M}|\xi(s)|,
 \qquad
 \|f\|_\cY
 =\sup_{J_\sigma}\e^{-c|s|}
   \langle s\rangle^{-(M-d_0)}|f(s)|.
\]
The Green operator is uniformly bounded from \(\cY\) to \(\cX\).
On the ball \(\|\xi\|_\cX\le C_0\del\), the nonlinear source satisfies
\begin{align}
 \|N(\cdot,\xi)\|_\cY
 &\le C\del+C\del^2\Lambda_R,\\
 \|N(\cdot,\xi)-N(\cdot,\widetilde\xi)\|_\cY
 &\le C\del\Lambda_R\|\xi-\widetilde\xi\|_\cX,
 \label{eq:appendix-N-lipschitz}
\end{align}
where
\[
 \Lambda_R=\e^{cR}\langle R\rangle^{M+m_0+d_0}.
\]
For example, a quadratic term is bounded by
\[
 \e^{-c|s|}\langle s\rangle^{-(M-d_0)}
 |\xi(s)|\,|\xi(s)-\widetilde\xi(s)|
 \le C\del\e^{cR}\langle R\rangle^{M+d_0}
       \|\xi-\widetilde\xi\|_\cX.
\]

At an endpoint, \cref{eq:nonlinear-boundary} and the implicit-function
theorem give
\[
 |\mathcal B(a)|\le CE(R)^{-1}|a|^2,
 \qquad
 |\mathcal B(a)-\mathcal B(\widetilde a)|
 \le CE(R)^{-1}(|a|+|\widetilde a|)|a-\widetilde a|.
\]
Multiplication by the normalized boundary weight produces the same
\(C\del\Lambda_R\) Lipschitz factor.  Because
\(R=S_\del+B\) and
\(\del\e^{cS_\del}\langle S_\del\rangle^M=o(1)\),
the combined nonlinear map is a strict contraction.  Parameter
differentiation gives a uniformly invertible linearized map; the direct
source orders \(O(\del)\) for \(\nu\) and \(O(\del^2)\) for \(\eta\)
then yield \cref{eq:trace-estimates}.

\subsection{Moving transition hits}
\label{app:moving-hits}

The auxiliary prepared endpoints are not physical transition sections.
As in the main text, write \(S=S_\del\) and use the sections in
\cref{eq:transition-sections}.  Let
\(t_a<0<t_r\) be the unique hit times of the two canonical traces and set
\(I_a=[t_a,0]\), \(I_r=[0,t_r]\), as in
\cref{eq:retained-intervals}.
Since
\[
 Y_0'(s)=\frac{s}{2\alpha},
\]
the trace estimates imply strict monotonicity on the unit neighborhoods of
\(\pm S_\del\), and hence
\begin{equation}
\label{eq:hitting-time-bounds}
 |t_a+S_\del|+|t_r-S_\del|
 \le C\del\e^{cS_\del}\langle S_\del\rangle^m.
\end{equation}
There is no second hit: on the outer pieces the sign of \(X\) fixes the
sign of \(Y'\), while the compact inner piece has bounded \(Y\) and the
section height diverges.

For a parameter \(\lambda\),
\[
 \partial_\lambda t
 =-\frac{\partial_\lambda Y(t)}{Y'(t)},
\]
and
\[
 \partial_{\eta\eta}t
 =-\frac{Y_{\eta\eta}+2Y_{s\eta}t_\eta+
                   Y_{ss}t_\eta^2}{Y_s}\bigg|_{s=t}.
\]
Using \cref{eq:trace-estimates} and
\(|Y'(t)|\ge S_\del/(4\alpha)\) gives the same bound as
\cref{eq:hitting-time-bounds} for a \(\nu\)-derivative, and the bound with
\(\del\) replaced by \(\del^2\) for every derivative containing
\(\eta\).  Mixed \(\nu\)--\(\eta\) derivatives obey the latter estimate.

\subsection{Gaussian suppression of the prepared boundary}

Suppose a phase-compatible current-state error at one endpoint, together
with its declared parameter derivatives, is bounded by
\(A\e^{cR}\langle R\rangle^m\).  Formula
\cref{eq:ab-coordinates} gives its normal coefficient as
\[
 |\beta|\le
 CA\e^{-R^2/2+cR}\langle R\rangle^{m+1}.
\]
At the central section \(n(0)=(0,1)^T\), and the tangent coefficient has
been removed by the common phase.  Therefore the contribution to the
central trace and its Gaussian integral norm is at most
\begin{equation}
\label{eq:boundary-suppression}
 CA\e^{-R^2/2+cR}\langle R\rangle^{m+2}.
\end{equation}
If two preparations agree on the retained tube but differ in the outer
annulus, the difference forcing is supported outside \(|s|\le S\).  The
tangent integral in \cref{eq:one-sided-green} vanishes at the central point,
and its normal part has the same bound with \(R\) replaced by \(S\).
The nonlinear resolvent changes only the constant.  Parameter
differentiation preserves the support and one-sided structure.

Because \(S_\del^2/2=p\log(1/\del)\), the right side of
\cref{eq:boundary-suppression} is
\[
 A\del^p\e^{cS_\del}\langle S_\del\rangle^{m+2}
 =A\del^{p-o(1)}.
\]
The tame exponents are fixed before \(p\) is chosen.  This is why the
quantifier order in \cref{thm:main} is necessary.

\subsection{Differentiating the gap}

On either trace,
\[
 \frac{\dd}{\dd s}\mathscr H(z)
 =\del A_1(z,\nu)+\del^2A_2(z,\nu,\eta)
  +\del^3A_3(z,\del,\nu,\eta),
\]
where \(A_k=\nabla\mathscr H\cdot Q_k\) for \(k=1,2\), and
\(A_3=\nabla\mathscr H\cdot R_{3,Q}\).  Splitting at zero and allowing
the moving transition hits gives
\begin{align}
 \frac{\alpha\e}{2}\mathfrak d_\cP
 ={}&\mathscr H(z^a(t_a))-\mathscr H(z^r(t_r))\notag\\
 &+\int_{t_a}^{0}\frac{\dd}{\dd s}\mathscr H(z^a(s))\,\dd s
  +\int_{0}^{t_r}\frac{\dd}{\dd s}\mathscr H(z^r(s))\,\dd s.
 \label{eq:moving-gap-identity}
\end{align}
For the canonical prepared tails, the first line is not literally zero at
the moving hits.  Using \(\mathscr H(z^a(-R))=
\mathscr H(z^r(R))=0\), it equals
\[
 \int_{-R}^{t_a}\frac{\dd}{\dd s}\mathscr H(z^a(s))\,\dd s
 +\int_{t_r}^{R}\frac{\dd}{\dd s}\mathscr H(z^r(s))\,\dd s.
\]
These are precisely the two prepared outer-annulus contributions, and
\cref{eq:boundary-suppression} controls them and all declared derivatives.
For a nonzero physical boundary mismatch the same estimate applies after
adding its phase-compatible endpoint error.  Terms from differentiating the
moving limits contain the same Gaussian endpoint factor and the hitting-time
jets above.

On the two integrals, the pointwise graph and trace estimates are dominated
by
\[
 C\e^{-s^2/2+c|s|}\langle s\rangle^m,
\]
which is integrable on \(\R\).  Dominated differentiation is therefore
uniform.  The leading \(\nu\)-source is
\[
 \frac{\alpha\e}{2}\del\e^{-s^2/2},
\]
while the leading \(\eta\)-source is
\[
 -\frac{\alpha\e}{2}c_\perp\del^2s^2\e^{-s^2/2}.
\]
Extending the integrals to the whole line costs
\(O(\del^{p-o(1)})\).  The Gaussian moments then give
\[
 \int_\R\e^{-s^2/2}\,\dd s
 =\int_\R s^2\e^{-s^2/2}\,\dd s=\sqrt{2\pi},
\]
which proves the two derivative coefficients in
\cref{eq:full-gap-expansion}.  Twice differentiating the integrand shows
that the direct \(Q_2\) source is \(O(\del^2)\), while every trace-derivative
product is smaller by \cref{eq:trace-estimates}.  This proves the
\(\eta\eta\)-bound.

For completeness, a parameter derivative of the lifted complete history is
not merely the derivative of the reduced trace.  It is
\[
 \partial_\lambda\varphi^\sigma
 =D\iota_{\del,\nu,\eta}(z^\sigma)
     \partial_\lambda z^\sigma
  +\partial_\lambda\iota_{\del,\nu,\eta}(z^\sigma),
 \qquad \lambda\in\{\nu,\eta\},
\]
with the ordinary Faà di Bruno terms at higher mixed order.  The
logarithmic graph estimates control both terms and their declared mixed
derivatives uniformly, which is the history-level estimate used in
\cref{prop:gap-expansion}.
